% !TEX program = xelatex
\documentclass{article}
\usepackage{geometry}\geometry{margin=23mm}
\usepackage{babel}\babelprovide[import,main]{ukrainian}\babelfont{rm}{Times New Roman}
\begin{document}
\title{Мірна дошка}\author{}\date{14 травня 2026 року}\maketitle
\begin{abstract}Дві позначки показують початкове положення дошки.\end{abstract}
\tableofcontents
\section{Позначки}\begin{quote}«Спершу прочитайте нижню позначку».\end{quote}
\begin{figure}[ht]\centering\fbox{\rule{0pt}{8mm}\rule{33mm}{0pt}}\caption{Положення дошки}\end{figure}
\section{Висота}
\begin{table}[ht]\centering\begin{tabular}{cr}Позначка&Висота (mm)\\\hline 1&12\\2&18\end{tabular}\caption{Висота позначок}\end{table}
\begin{thebibliography}{1}\bibitem{record}\emph{Зошит вимірювань}, 2025, \chaptername~2.\end{thebibliography}
\end{document}
